% !PRINTED 310
\weberpage{310}
%<<CONT
$x = -\infty$ alle Zeichenfolgen verloren gegangen, also in der Doppelreihe der $f_\nu, F_\nu$ für $x = +\infty$ alle $VP$, für $x = -\infty$ alle $PP$.

Danach können wir nach \S.\,100, II, III folgende beiden Theoreme aussprechen:

\gesperrt{V. Die Zahl der positiven Wurzeln von $f(x)$ ist so gross oder um eine gerade Zahl kleiner als die Zahl der Wechsel-Folgen in der Doppelreihe (4).}

\gesperrt{VI. Die Zahl der negativen Wurzeln von $f(x)$ ist so gross oder um eine gerade Zahl kleiner als die Zahl der Folge-Folgen in der Doppelreihe (4).}

Die Summe der Anzahlen der $VP$ und der $PP$ ist aber gleich der Anzahl der Zeichenfolgen in der einfachen Reihe
\weq{5}{A_0, A_1, A_2 \ldots A_n,}
und wir können also V. und VI. in den einen Satz zusammenfassen:

\gesperrt{VII. Die Anzahl aller reeller Wurzeln von $f(x)$ ist so gross oder um eine gerade Zahl kleiner als die Zahl der Zeichenfolgen in der Reihe (5).}

Die Gesammtanzahl der Zeichenwechsel und Zeichenfolgen in der Reihe (5) ist so gross, wie die Zahl aller reellen und imaginären Wurzeln zusammengenommen, nämlich $n$. Denn der Fall, dass in der Reihe (5) einige Glieder verschwinden, ist hier auszunehmen. Da überdies das erste und letzte Glied der Reihe (5) dasselbe Vorzeichen haben, so ist die Anzahl der Zeichenwechsel in dieser Reihe eine gerade Zahl. Wir können also zu VII., ohne damit etwas Neues zu sagen, noch das Theorem hinzufügen:

\gesperrt{VIII. Die Anzahl der imaginären Wurzeln von $f(x)$ ist mindestens gleich der Anzahl der Zeichenwechsel in der Reihe (5).}

In der Fassung der beiden letzten Sätze ist es gleichgültig, ob einige der Coëfficienten $a_1, a_2 \ldots a_n$ verschwinden, wenn nur keines der $A_\nu$ Null ist.

Bei der cubischen Gleichung
\[
a_0 x^3 + a_1 x^2 + a_2 x + a_3 = 0
\]
ist
\[
A_1 = 2(a_1^2 - 3a_0a_2)
\]
\[
A_2 = 2(a_2^2 - 3a_1a_3),
\]
%CONT>>
